\subsection{Central Limit Theory}
% Characteristic functions, CLT, Lindeberg/Feller flavor, delta method preview, weak convergence.

The law of large numbers says that $\overline{X}_n$ is close to $\mu$ for large
$n$. The central limit theorem describes the scale and shape of the remaining
fluctuation.

\begin{theorem} \label{thm:central-limit-theorem-iid} \hlt{Central Limit Theorem}\\
Let $(X_i)_{i\geq1}$ be IID real-valued random variables, $\mathbb{E}X_1=\mu$, $\operatorname{Var}(X_1)=\sigma^2\in(0,+\infty)$.\\
Let $\overline{X}_n\defeq n^{-1}\sum_{i=1}^{n}X_i$, $Z_n\defeq\sigma^{-1}\sqrt{n}(\overline{X}_n-\mu)$.\\
Then $Z_n\rightsquigarrow N(0,1)$, equivalently $\mathbb{P}(Z_n\leq z)\to\Phi(z)\defeq\int_{-\infty}^{z}(2\pi)^{-1/2}e^{-x^2/2}\,dx$ $\forall z\in\mathbb{R}$.
\begin{proof} Set $Y_i\defeq(X_i-\mu)/\sigma$, so $\mathbb{E}Y_i=0$, $\operatorname{Var}(Y_i)=1$, and $Z_n=n^{-1/2}\sum_{i=1}^{n}Y_i$.\\
For $\varphi(t)\defeq\mathbb{E}e^{\mathrm{i}tY_1}$, the finite second moment expansion gives $\varphi(u)=1-u^2/2+o(u^2)$ as $u\to0$. \\
By independence, $\varphi_{Z_n}(t)=[\varphi(t/\sqrt{n})]^n=[1-t^2/(2n)+o(1/n)]^n\to e^{-t^2/2}$ for every $t\in\mathbb{R}$. \\
Since $e^{-t^2/2}$ is the characteristic function of $N(0,1)$, Levy's continuity theorem gives $Z_n\rightsquigarrow N(0,1)$; continuity of $\Phi$ gives the displayed c.d.f. convergence.\\
\hlt{Structural Commentary}: Centering and scaling turn the sum into unit-variance IID increments.\\
\hlt{Key Insight}: Characteristic functions turn normalized sums into powers, whose limit is Gaussian.
\end{proof}
\end{theorem}

\begin{method} \label{meth:clt-normal-approximation-forms} \hlt{Equivalent CLT Approximation Forms}\\
The CLT approximation $Z_n\rightsquigarrow N(0,1)$ is used as $\overline{X}_n\approx N(\mu,\sigma^2/n)$, equivalently $\sqrt{n}(\overline{X}_n-\mu)\approx N(0,\sigma^2)$, for approximating probabilities involving $\overline{X}_n$.
\end{method}

\begin{lemma} \label{lem:mgf-convergence-criterion} \hlt{MGF Convergence Criterion}\\
Let $(Z_n)_{n\geq1}$ and $Z$ be real-valued random variables with MGFs $\psi_n(t)\defeq\mathbb{E}e^{tZ_n}$ and $\psi(t)\defeq\mathbb{E}e^{tZ}$ finite on $(-a,a)$ for some $a>0$. If $\psi_n(t)\to\psi(t)$ for every $t\in(-a,a)$, then $Z_n\rightsquigarrow Z$.
\begin{proof} Assume that $Y_i\defeq(X_i-\mu)/\sigma$ has MGF $\psi$ finite near $0$.\\
Then $\mathbb{E}Y_i=0$, $\operatorname{Var}(Y_i)=1$, and $\psi(t)=1+t^2/2+o(t^2)$ as $t\to0$. Since $Z_n=n^{-1/2}\sum_{i=1}^{n}Y_i$,\\
independence gives $\psi_{Z_n}(t)=[\psi(t/\sqrt{n})]^n=[1+t^2/(2n)+o(1/n)]^n\to e^{t^2/2}$ for every $t$ near $0$.\\
As $e^{t^2/2}$ is the MGF of $N(0,1)$, Lemma \ref{lem:mgf-convergence-criterion} gives $Z_n\rightsquigarrow N(0,1)$.\\
\hlt{Structural Commentary}: Standardization reduces the MGF to $1+t^2/2+o(t^2)$.\\
\hlt{Key Insight}: Independence turns normalized sums into MGF powers converging to the Gaussian MGF.
\end{proof}
\end{lemma}

\begin{example} \label{ex:clt-poisson-program-errors} \hlt{A Poisson Error-Count Approximation}\\
Let $X_1,\ldots,X_{125}\sim\operatorname{Poisson}(5)$ be IID and set $\overline{X}_{125}=125^{-1}\sum_{i=1}^{125}X_i$.\\
Since $\mu=5$ and $\sigma^2=5$, the CLT gives $\mathbb{P}(\overline{X}_{125}<5.5)\approx\Phi(\sqrt{125}(5.5-5)/\sqrt5)=\Phi(2.5)\approx0.9938$.
\end{example}

\begin{theorem} \label{thm:studentized-central-limit-theorem} \hlt{Studentized Central Limit Theorem}\\
Let $(X_i)_{i\geq1}$ be IID real-valued random variables with $\mathbb{E}X_1=\mu$ and $\operatorname{Var}(X_1)=\sigma^2\in(0,+\infty)$.\\
Set $\overline{X}_n\defeq n^{-1}\sum_{i=1}^{n}X_i$, $S_n^2\defeq(n-1)^{-1}\sum_{i=1}^{n}(X_i-\overline{X}_n)^2$, and $T_n\defeq S_n^{-1}\sqrt{n}(\overline{X}_n-\mu)$.
Then $T_n\rightsquigarrow N(0,1)$.
\begin{proof}
The CLT gives $Z_n\defeq\sigma^{-1}\sqrt{n}(\overline{X}_n-\mu)\rightsquigarrow N(0,1)$, and Proposition \ref{prop:sample-variance-consistency} gives $S_n^2\xrightarrow{\mathbb{P}}\sigma^2$, hence $S_n\xrightarrow{\mathbb{P}}\sigma$.\\
Thus $\sigma/S_n\xrightarrow{\mathbb{P}}1$, and Slutsky's rule from Theorem \ref{thm:continuous-mapping-and-slutsky} gives $T_n=Z_n(\sigma/S_n)\rightsquigarrow N(0,1)$.\\
\hlt{Structural Commentary}: Studentization replaces the unknown $\sigma$ by the consistent estimate $S_n$.\\
\hlt{Key Insight}: Slutsky permits multiplying the CLT statistic by $\sigma/S_n\to1$.
\end{proof}
\end{theorem}

\begin{theorem} \label{thm:berry-esseen-rate} \hlt{Berry--Esseen Rate}\\
Let $(X_i)_{i\geq1}$ be IID real-valued random variables, $\mathbb{E}X_1=\mu$, $\operatorname{Var}(X_1)=\sigma^2\in(0,+\infty)$, $\rho_3\defeq\mathbb{E}\abs{X_1-\mu}^3<+\infty$.\\
Let $\overline{X}_n\defeq n^{-1}\sum_{i=1}^{n}X_i$, $Z_n\defeq\sigma^{-1}\sqrt{n}(\overline{X}_n-\mu)$, and $\Phi(z)\defeq\int_{-\infty}^{z}(2\pi)^{-1/2}e^{-x^2/2}\,dx$.\\
Then $\sup_{z\in\mathbb{R}}\abs{\mathbb{P}(Z_n\leq z)-\Phi(z)}\leq(33/4)\rho_3/(\sigma^3\sqrt{n})$.
\begin{proof} Set $Y_i\defeq(X_i-\mu)/\sigma$, $\rho\defeq\mathbb{E}\abs{Y_1}^3=\rho_3/\sigma^3$, and $W_n\defeq n^{-1/2}\sum_{i=1}^{n}Y_i=Z_n$.\\
For $\varphi(t)\defeq\mathbb{E}e^{\mathrm{i}tY_1}$, Taylor's formula with $\mathbb{E}Y_1=0$, $\mathbb{E}Y_1^2=1$, and $\mathbb{E}\abs{Y_1}^3=\rho$ gives $\varphi(t)=1-t^2/2+O(\rho\abs{t}^3)$ near $0$. Thus $\varphi_{W_n}(t)=[\varphi(t/\sqrt{n})]^n$ differs from $e^{-t^2/2}$ by $O(\rho n^{-1/2})$ in Esseen's smoothing integral, and the standard optimized smoothing bound gives $\sup_z\abs{\mathbb{P}(W_n\leq z)-\Phi(z)}\leq(33/4)\rho/\sqrt{n}$.\\
Substituting $\rho=\rho_3/\sigma^3$ gives the stated inequality.\\
\hlt{Structural Commentary}: Berry--Esseen bounds the CLT error by standardized third absolute moment.\\
\hlt{Key Insight}: Normalize, compare characteristic functions, then apply Esseen smoothing.
\end{proof}
\end{theorem}

\begin{theorem} \label{thm:multivariate-central-limit-theorem} \hlt{Multivariate Central Limit Theorem}\\
Let $(X_i)_{i\geq1}$ be IID random vectors in $\mathbb{R}^k$ with $\mathbb{E}X_1=\mu\in\mathbb{R}^k$ and $\operatorname{Cov}(X_1)=\Sigma$, and set $\overline{X}_n\defeq n^{-1}\sum_{i=1}^{n}X_i$.\\
Then $\sqrt{n}(\overline{X}_n-\mu)\rightsquigarrow N_k(0,\Sigma)$.
\begin{proof}
For any $a\in\mathbb{R}^k$, the variables $a^tX_i$ are IID with mean $a^t\mu$ and variance $a^t\Sigma a$.\\
If $a^t\Sigma a>0$, the univariate CLT gives $a^t\sqrt{n}(\overline{X}_n-\mu)\rightsquigarrow N(0,a^t\Sigma a)$; if $a^t\Sigma a=0$, the projection is identically $0$ a.s. after centering. Thus every linear projection converges to the corresponding projection of $N_k(0,\Sigma)$, and Cramer--Wold gives $\sqrt{n}(\overline{X}_n-\mu)\rightsquigarrow N_k(0,\Sigma)$.\\
\hlt{Structural Commentary}: Multivariate convergence is checked through one-dimensional projections.\\
\hlt{Key Insight}: Project first, apply the scalar CLT, then use Cramer--Wold.
\end{proof}
\end{theorem}

\begin{theorem} \label{thm:delta-method-univariate} \hlt{Delta Method}\\
Let $(Y_n)_{n\geq1}$ be real-valued random variables, let $\mu\in\mathbb{R}$ and $\sigma>0$, and assume $\sigma^{-1}\sqrt{n}(Y_n-\mu)\rightsquigarrow N(0,1)$.\\
If $g:\mathbb{R}\to\mathbb{R}$ is differentiable at $\mu$ and $g'(\mu)\neq0$, then $(\abs{g'(\mu)}\sigma)^{-1}\sqrt{n}(g(Y_n)-g(\mu))\rightsquigarrow N(0,1)$.\\
Equivalently $g(Y_n)\approx N(g(\mu),(g'(\mu))^2\sigma^2/n)$.
\begin{proof} Set $A_n\defeq\sigma^{-1}\sqrt{n}(Y_n-\mu)$, so $A_n\rightsquigarrow N(0,1)$ and $Y_n\xrightarrow{\mathbb{P}}\mu$.\\
Differentiability gives $g(Y_n)-g(\mu)=g'(\mu)(Y_n-\mu)+r(Y_n)(Y_n-\mu)$, where $r(y)\to0$ as $y\to\mu$, hence $r(Y_n)\xrightarrow{\mathbb{P}}0$.\\
Therefore $(\abs{g'(\mu)}\sigma)^{-1}\sqrt{n}(g(Y_n)-g(\mu))=\operatorname{sgn}(g'(\mu))A_n(1+r(Y_n)/g'(\mu))\rightsquigarrow N(0,1)$ by Theorem \ref{thm:continuous-mapping-and-slutsky} and symmetry of the standard normal.\\
\hlt{Structural Commentary}: The delta method is first-order Taylor expansion at the limit point.\\
\hlt{Key Insight}: The derivative rescales the fluctuation; the remainder is negligible in probability.
\end{proof}
\end{theorem}

\begin{example} \label{ex:delta-method-exponential-sample-mean} \hlt{Exponentiating a Sample Mean}\\
Let $(X_i)_{i=1}^{n}$ be IID with mean $\mu$ and variance $\sigma^2<+\infty$, and set $W_n\defeq e^{\overline{X}_n}$.\\
For $g(s)=e^s$, $g'(\mu)=e^\mu$, so the delta method gives $W_n\approx N(e^\mu,e^{2\mu}\sigma^2/n)$.
\end{example}

\begin{theorem} \label{thm:delta-method-multivariate} \hlt{Multivariate Delta Method}\\
Let $(Y_n)_{n\geq1}$ be random vectors in $\mathbb{R}^k$ with $\sqrt{n}(Y_n-\mu)\rightsquigarrow N_k(0,\Sigma)$, and let $g:\mathbb{R}^k\to\mathbb{R}$ be differentiable at $\mu$ with gradient $\nabla g(\mu)$. Then $\sqrt{n}(g(Y_n)-g(\mu))\rightsquigarrow N(0,\nabla g(\mu)^t\Sigma\nabla g(\mu))$.
\begin{proof} Set $A_n\defeq\sqrt{n}(Y_n-\mu)$, so $A_n\rightsquigarrow N_k(0,\Sigma)$ and $Y_n\xrightarrow{\mathbb{P}}\mu$.\\
Differentiability gives $g(Y_n)-g(\mu)=\nabla g(\mu)^t(Y_n-\mu)+r(Y_n)\norm{Y_n-\mu}$ with $r(Y_n)\xrightarrow{\mathbb{P}}0$.\\
Since $A_n=O_{\mathbb{P}}(1)$, $\sqrt{n}r(Y_n)\norm{Y_n-\mu}=r(Y_n)\norm{A_n}\xrightarrow{\mathbb{P}}0$.\\
Thus $\sqrt{n}(g(Y_n)-g(\mu))=\nabla g(\mu)^tA_n+o_{\mathbb{P}}(1)\rightsquigarrow N(0,\nabla g(\mu)^t\Sigma\nabla g(\mu))$ by Theorem \ref{thm:continuous-mapping-and-slutsky}.\\
\hlt{Structural Commentary}: Multivariate delta method linearizes $g$ by its gradient.\\
\hlt{Key Insight}: The gradient projection carries the Gaussian limit; the differentiability remainder vanishes.
\end{proof}
\end{theorem}

\begin{example} \label{ex:multivariate-delta-method-product-and-ratio} \hlt{Product and Ratio of Sample Means}\\
Let $X_i=(X_{1i},X_{2i})^t$ be IID with mean $\mu=(\mu_1,\mu_2)^t$ and covariance $\Sigma=(\sigma_{ij})$, and set $\overline{X}_{jn}\defeq n^{-1}\sum_{i=1}^{n}X_{ji}$.\\
For $g(x_1,x_2)=x_1x_2$, $\nabla g(\mu)=(\mu_2,\mu_1)^t$, so $\sqrt{n}(\overline{X}_{1n}\overline{X}_{2n}-\mu_1\mu_2)\rightsquigarrow N(0,\mu_2^2\sigma_{11}+2\mu_1\mu_2\sigma_{12}+\mu_1^2\sigma_{22})$.\\
If $\mu_2\neq0$, then for $h(x_1,x_2)=x_1/x_2$ and $a\defeq\nabla h(\mu)=(1/\mu_2,-\mu_1/\mu_2^2)^t$, $\sqrt{n}(\overline{X}_{1n}/\overline{X}_{2n}-\mu_1/\mu_2)\rightsquigarrow N(0,a^t\Sigma a)$.
\end{example}
